\chapter{Geltungsbereich und Grundbegriffe}

Dieser Leitfaden sammelt zentrale Resultate, bei denen Endlichkeit
eine Existenz, eine Stabilisierung, eine einheitliche endliche Schranke
oder ein endliches Verfahren sichert. Allgemeingültige Sätze,
abtrennbare Teilaussagen und Richtungen werden nicht mitgeführt.
Die Beschränkung eines Begriffs auf endliche Halbgruppen allein ist
kein Aufnahmegrund.
Benötigte allgemeine Begriffe werden weiterhin definiert.

Die Kapitel bilden zugleich eine mögliche Gliederung für den weiteren
Ausbau von Band~37: Potenzen, innere Struktur, Ideale, endliche
Kanonisierung, Normalformen, Einbettungen und Überlagerungen, globale Zerlegung und
relationale Abbildungen.
Der Leitfaden enthält Definitionen und Hauptsätze, keine Beweise und
keine Beispiele. Er ist eine Auswahl zentraler Resultate, keine
abgeschlossene Klassifikation aller endlichen Halbgruppen.

An die betrachtete Halbgruppe werden nur Endlichkeit und Assoziativität
gestellt. Ein neutrales Element, ein Nullelement, Kommutativität,
Kürzbarkeit, Regularität oder Aperiodizität werden nicht vorausgesetzt.
Interne Gruppen, lokale Einheiten und Rees-Konstruktionen treten als
\emph{Ergebnisse} der Strukturuntersuchung oder als Darstellungsbausteine auf.
Die Aussagen über ausgewählte Abbildungen, Unterhalbgruppen oder Klassen
benennen zusätzlich die jeweils betrachteten Objekte. Bedingungen an
Zielhalbgruppen relationaler Morphismen sind keine Zusatzannahmen an $S$.

Wie im Hauptprojekt sind leere Halbgruppen zugelassen. Nur wenn ein Satz
die Existenz eines Elements oder eines nichtleeren Ideals behauptet,
wird der nichtleere Fall ausdrücklich genannt. Endlichkeit muß nicht
die schwächste mögliche Voraussetzung eines aufgenommenen Satzes sein:
Er darf auch in besonderen unendlichen Halbgruppen gelten, jedoch
nicht schon in beliebigen Halbgruppen ohne weitere Voraussetzungen.

\section{Gemeinsame Sprache}

\begin{definition}[Endliche Halbgruppe und Grundnotation]
Eine \emph{endliche Halbgruppe} ist eine endliche Menge $S$ mit einer
Abbildung $S\times S\to S$, $(a,b)\mapsto ab$, für die
\[
 (ab)c=a(bc)\qquad(a,b,c\in S)
\]
gilt. Im ganzen Leitfaden bezeichnet $S$ eine solche Halbgruppe und
$n=|S|$ ihre Größe. Es sei $\mathbb N_{>0}=\{1,2,\ldots\}$.
Für $a\in S$ definieren wir $a^1=a$ und $a^{k+1}=a^k a$.
Für $A,B\subseteq S$ und $k\geq1$ setzen wir
\[
 AB=\{ab:a\in A,\ b\in B\},\qquad
 S^k=\{s_1\cdots s_k:s_1,\ldots,s_k\in S\}.
\]
Eine nichtleere produktabgeschlossene Teilmenge heißt
\emph{Unterhalbgruppe}; insbesondere
$\langle a\rangle=\{a^k:k\geq1\}$.
\end{definition}

\begin{definition}[Idempotente und adjungierte Einheit]
Ein Element $e\in S$ heißt \emph{idempotent}, wenn $e^2=e$.
Es sei $E(S)=\{e\in S:e^2=e\}$.
Besitzt $S$ bereits ein zweiseitiges neutrales Element, setzen wir
$S^1=S$. Andernfalls ist $S^1=S\mathbin{\dot\cup}\{1\}$ mit
$1x=x1=x$ für $x\in S^1$. Diese Notation setzt kein neutrales Element
in $S$ voraus; positive Potenzen werden weiterhin in $S$ gebildet.
\end{definition}

\begin{definition}[Strukturerhaltende Abbildungen]
Für Halbgruppen $S,T$ heißt $f:S\to T$ \emph{Homomorphismus}, wenn
$f(ab)=f(a)f(b)$ für alle $a,b\in S$ gilt. Ein bijektiver
Homomorphismus heißt \emph{Isomorphismus}; $S\cong T$ bedeutet,
dass ein solcher existiert.
\end{definition}
